\documentclass[12pt,a4paper]{article}
\usepackage[utf8]{inputenc}
\usepackage[T2A]{fontenc}
\usepackage[russian]{babel}
\usepackage{amsmath,amssymb}
\usepackage{enumitem}

\usepackage[left=25mm,right=15mm,top=20mm,bottom=20mm]{geometry}

\title{Ограничения применимости псевдовремени\\
в однофазной и многофазной фильтрации}
\author{}
\date{}

\begin{document}
\maketitle

%=============================================================================
\section{Однофазный случай (реальный газ)}\label{sec:1ph}
%=============================================================================

\subsection{Псевдодавление: точная замена}\label{ssec:mp_exact}

Уравнение фильтрации реального газа:
\begin{equation}\label{eq:gas_filt}
  \phi\,\frac{\partial}{\partial t}\left(\frac{p}{Z}\right)
  = \mathrm{div}\left(\frac{kp}{\mu_g Z}\,\nabla p\right).
\end{equation}
Псевдодавление Аль-Хуссейни (1966):
\begin{equation}\label{eq:mp}
  m(p) = 2\int_{p_{\mathrm{ref}}}^{p}\frac{p'\,dp'}{\mu_g(p')\,Z(p')}.
\end{equation}
Подстановка~\eqref{eq:mp} в~\eqref{eq:gas_filt} \emph{точно}
устраняет нелинейность пространственного оператора:
\begin{equation}\label{eq:gas_m}
  \phi\,\mu_g(p)\,c_g(p)\,\frac{\partial m}{\partial t}
  = k\,\nabla^2 m.
\end{equation}
Здесь $c_g = (1/p - (1/Z)\,dZ/dp)$ --- сжимаемость газа.
Замена $p \to m$ \textbf{не является приближением}
и справедлива при любых перепадах давления.

Оставшаяся нелинейность~--- в коэффициенте $\mu_g c_g$ перед
$\partial m/\partial t$. Для устранения и этой нелинейности
Agarwal~(1979) вводит псевдовремя.

\subsection{Псевдовремя: точное определение}\label{ssec:pt_exact}

Точное определение:
\begin{equation}\label{eq:ta_exact}
  t_a(r,\,t) = \int_0^t
  \frac{\mu_i\,c_{gi}}{\mu_g\bigl(p(r,\tau)\bigr)\,
        c_g\bigl(p(r,\tau)\bigr)}\,d\tau,
\end{equation}
где $\mu_i = \mu_g(p_0)$, $c_{gi} = c_g(p_0)$.

$t_a$ определено как \emph{пространственно-временн\'{о}е поле},
не как функция одного аргумента~$t$:
в каждой точке~$r$ давление проходит свою историю $p(r,\tau)$,
и подинтегральное выражение зависит от~$r$.

\subsection{Деформация пространственного оператора}
\label{ssec:deformation}

Замена координат $(r,t) \to (r,t_a)$, где $t_a = t_a(r,t)$,
деформирует пространственный оператор.

\paragraph{Производная по~$r$ при фиксированном~$t_a$.}
\begin{equation}\label{eq:dr_ta}
  \left.\frac{\partial}{\partial r}\right|_{t_a}
  = \left.\frac{\partial}{\partial r}\right|_{t}
  + \left.\frac{\partial t}{\partial r}\right|_{t_a}
    \frac{\partial}{\partial t}\bigg|_{r}.
\end{equation}
Множитель $(\partial t/\partial r)|_{t_a}$ отличен от нуля
всюду, где давление неоднородно. Из~\eqref{eq:ta_exact}:
\begin{equation}\label{eq:dta_dr}
  \frac{\partial t_a}{\partial r}
  = \int_0^t \mu_i c_{gi}\,
  \frac{\partial}{\partial r}
  \!\left[\frac{1}{\mu_g\bigl(p(r,\tau)\bigr)\,
  c_g\bigl(p(r,\tau)\bigr)}\right] d\tau
  \neq 0.
\end{equation}

\paragraph{Деформированный оператор.}
Лапласиан $\nabla^2 m$ в координатах $(r,t_a)$
приобретает дополнительные члены. Раскрывая
$\partial^2 m/\partial r^2|_{t_a}$ с помощью~\eqref{eq:dr_ta}:
\begin{equation}\label{eq:laplacian_deformed}
  \nabla^2 m\Big|_{t} = \nabla^2 m\Big|_{t_a}
  - 2\,\frac{\partial t}{\partial r}\bigg|_{t_a}
    \frac{\partial^2 m}{\partial r\,\partial t_a}
  - \left[\frac{\partial^2 t}{\partial r^2}\bigg|_{t_a}
  + \frac{1}{r}\,\frac{\partial t}{\partial r}\bigg|_{t_a}\right]
    \frac{\partial m}{\partial t_a}
  + \ldots
\end{equation}
(в цилиндрических координатах; $\ldots$ содержит
$(\partial t/\partial r)^2\,\partial^2 m/\partial t_a^2$).

\paragraph{Вывод.}
В координатах $(r,t_a)$ уравнение диффузии
$\phi\mu_ic_{gi}\,\partial m/\partial t_a = k\nabla^2 m$
\textbf{не выполняется}: лапласиан содержит перекрёстный
член $\partial^2 m/\partial r\,\partial t_a$ и
конвективный член $\partial m/\partial t_a$.
Уравнение перестаёт быть уравнением диффузии.

Осреднение $p(r,\tau) \to \bar{p}(\tau)$~--- это
\emph{не аппроксимация точного определения}~\eqref{eq:ta_exact},
а единственный способ сохранить структуру уравнения.

\subsection{Стандартное приближение: условия применимости}
\label{ssec:conditions_1ph}

Заменяем $p(r,\tau) \to \bar{p}(\tau)$:
\begin{equation}\label{eq:ta_std}
  t_a(t) = \int_0^t
  \frac{\mu_i\,c_{gi}}{\mu_g\bigl(\bar{p}(\tau)\bigr)\,
        c_g\bigl(\bar{p}(\tau)\bigr)}\,d\tau.
\end{equation}
Теперь $t_a$~--- функция одного аргумента,
$\partial t_a/\partial r = 0$, и лапласиан
не деформируется. Уравнение:
\begin{equation}\label{eq:gas_linear}
  \phi\,\mu_i\,c_{gi}\,\frac{\partial m}{\partial t_a}
  = k\,\nabla^2 m
\end{equation}
формально линейно (все коэффициенты --- константы).

\paragraph{Физический смысл приближения.}
Замена $\mu_g c_g(p(r,\tau))$ на $\mu_g c_g(\bar{p}(\tau))$
корректна, когда пространственные вариации $p(r,\tau)$
вокруг среднего $\bar{p}(\tau)$ \emph{не приводят к существенным
изменениям} произведения $\mu_g c_g$. Формализуем это.

Разложим подинтегральное выражение~\eqref{eq:ta_exact}
в ряд по отклонению $\delta p(r,\tau) = p(r,\tau) - \bar{p}(\tau)$:
\begin{equation}\label{eq:taylor_mu_c}
  \frac{1}{\mu_g c_g(p)} = \frac{1}{\mu_g c_g(\bar{p})}
  + \frac{d}{dp}\left(\frac{1}{\mu_g c_g}\right)\bigg|_{\bar{p}}
  \!\!\delta p + O(\delta p^2).
\end{equation}
Приближение $t_a(r,t) \approx t_a(t)$ означает отбрасывание
членов $O(\delta p)$ и выше. Это корректно при условии:
\begin{equation}\label{eq:criterion_1ph}
  \boxed{
  \left|\frac{d\ln(\mu_g c_g)}{dp}\right| \cdot |\delta p|
  \ll 1,}
\end{equation}
где $\delta p = p - \bar{p}$~--- отклонение давления
от среднего в любой точке пласта.

\paragraph{Два масштаба, определяющие применимость.}

\begin{enumerate}
  \item \textbf{Масштаб неоднородности давления}
        $\delta p = p_{\max} - p_{\min}$.
        В радиальном потоке основной перепад давления
        приходится на ПЗП (несколько радиусов скважины),
        тогда как вдали от скважины давление почти однородно.
        При квазистационаре: $\delta p \sim (p_0 - p_{wf})$
        вблизи скважины, $\delta p \ll (p_0 - p_{wf})$
        в основном объёме пласта.

  \item \textbf{Масштаб нелинейности $\mu_g c_g$}~---
        характерное давление $p^*$, на котором $\mu_g c_g$
        изменяется в $e$~раз:
        \begin{equation}\label{eq:p_star}
          p^* = \left|\frac{d\ln(\mu_g c_g)}{dp}
          \right|^{-1}.
        \end{equation}
        Для идеального газа $c_g \sim 1/p$, $\mu_g \approx \mathrm{const}$,
        $p^* \sim p$. Для реального газа вблизи давления
        насыщения $p^*$ может быть значительно меньше.
\end{enumerate}

\noindent
\textbf{Критерий~\eqref{eq:criterion_1ph} в терминах масштабов:}
\begin{equation}\label{eq:criterion_scales}
  \boxed{\frac{\delta p}{p^*} \ll 1.}
\end{equation}

\paragraph{Когда критерий нарушается.}
\begin{enumerate}[nosep]
  \item \textbf{Большие депрессии.}
        При $\Delta p = p_0 - p_{wf} \sim p_0/2$ и $p^* \sim p_0$:
        $\delta p / p^* \sim 1/2$~--- на пределе.
  \item \textbf{Низкие давления.}
        При $p \to 0$ масштаб $p^* \to 0$ (для идеального газа
        $p^* = p$), и любой конечный $\delta p$ нарушает критерий.
  \item \textbf{Давление вблизи $p_b$.}
        Резкое изменение $c_g$ при фазовом переходе приводит
        к малому $p^*$ и нарушению~\eqref{eq:criterion_scales}.
  \item \textbf{Транзиентный поток.}
        При раннем транзиенте возмущение давления локализовано
        вблизи скважины, а среднее $\bar{p}$ определяется
        \emph{всем} объёмом пласта (включая невозмущённую часть).
        Отклонение $\delta p$ в ПЗП может быть сравнимо
        с полной депрессией.
\end{enumerate}

\paragraph{Выбор $\bar{p}$: два подхода.}

\begin{enumerate}
  \item \textbf{Agarwal (1979): $\bar{p}$ из материального баланса.}
        Среднее давление определяется по всему пласту:
        $\bar{p}/\bar{Z} = (p_i/Z_i)(1 - G_p/G)$.
        Корректно при boundary-dominated flow (BDF),
        когда весь объём участвует в дренировании.
        При транзиентном потоке $\bar{p}$ нечувствительно
        к фактическому фронту~--- возникает артефакт
        зависимости от далёких границ.

  \item \textbf{Anderson \& Mattar (2005): $\tilde{p}$ из области влияния.}
        Среднее давление определяется внутри области
        $A_r = \pi r_{\mathrm{inv}}^2 \cap A$:
        $\tilde{p}/\tilde{Z} = (p_i/Z_i)(1 - G_p/G_r)$,
        где $G_r$~--- запасы газа в области влияния.
        При переходе к BDF: $A_r \to A$, $\tilde{p} \to \bar{p}$.
        Этот подход минимизирует $\delta p$ в зоне,
        где $\mu_g c_g$ фактически <<работает>>, и
        расширяет область применимости псевдовремени
        на транзиентный режим.
\end{enumerate}

\paragraph{Замечание.}
Оба подхода предполагают, что профиль давления
\emph{внутри} области осреднения достаточно однородный.
Для ПЗП (призабойной зоны пласта) это условие
может нарушаться: давление падает логарифмически,
и перепад в пределах нескольких $r_w$ сравним
с полной депрессией. Формальная линейность~\eqref{eq:gas_linear}
не гарантирует физическую точность в ПЗП.


%=============================================================================
\section{Многофазный изотермический случай}\label{sec:multi}
%=============================================================================

\subsection{Новые степени свободы}

При переходе к системе Black~Oil (нефть с растворённым газом +
свободный газ) появляются три дополнительных источника сложности:

\begin{enumerate}
  \item \textbf{Насыщенность~--- второе неизвестное поле.}
        В однофазном случае состояние системы определяется
        одним полем $p(r,t)$. В двухфазном~--- двумя:
        $p(r,t)$ и $S_o(r,t)$, связанными системой уравнений.

  \item \textbf{Межфазный массообмен.}
        Газосодержание $R_s(p)$ связывает фазы:
        при $p < p_b$ газ выделяется из нефти, порождая
        члены $R_s'(dS_o/dp)$ в общей сжимаемости~$C_t$.

  \item \textbf{Относительные фазовые проницаемости.}
        Подвижности $k_{ro}(S_o)/\mu_o$, $k_{rg}(S_g)/\mu_g$
        зависят от насыщенности~--- нелинейность,
        отсутствующая в однофазной задаче.
\end{enumerate}

\subsection{Псевдодавление $m_g$: функция двух переменных}

\paragraph{Определение.}
Уравнение неразрывности газового компонента (Black~Oil, $R_v=0$):
\begin{equation}\label{eq:gas_cont}
  \frac{\partial}{\partial t}\left[\phi\left(\frac{S_g}{B_g}
  + R_s\,\frac{S_o}{B_o}\right)\right]
  - \mathrm{div}\left[k\left(\frac{k_{rg}}{B_g\mu_g}
  + R_s\,\frac{k_{ro}}{B_o\mu_o}\right)\nabla p\right] = 0.
\end{equation}
По аналогии с однофазным случаем, определим функцию
псевдодавления через суммарную подвижность газового
компонента. Введём обозначение:
\begin{equation}\label{eq:alpha_def}
  \alpha_g(S_o,\,p) = \frac{k_{rg}(1-S_o)}{B_g(p)\,\mu_g(p)}
  + R_s(p)\,\frac{k_{ro}(S_o)}{B_o(p)\,\mu_o(p)}
\end{equation}
--- суммарная подвижность газового компонента в обеих фазах.

В однофазной задаче подынтегральное выражение
$p/(\mu_g Z)$ определяется только давлением,
и $m(p)$~--- функция одного аргумента.
В~двухфазной задаче подвижность $\alpha_g$ зависит
от \emph{двух независимых переменных}: давления~$p$
(через PVT-свойства $B_g$, $\mu_g$, $R_s$, $B_o$, $\mu_o$)
и насыщенности~$S_o$ (через ОФП $k_{rg}$, $k_{ro}$).
Насыщенность~$S_o$~--- свободная переменная, не связанная
с давлением: она определяется начальными и граничными
условиями, историей вытеснения, ОФП.

Определим $m_g$ как параметрический интеграл,
в котором $S_o$~--- \textbf{свободный параметр},
а интегрирование ведётся по давлению:
\begin{equation}\label{eq:mg_multi}
  \boxed{m_g(p,\,S_o) = \mu_{gi}\,B_{gi}\int_{p_{\mathrm{ref}}}^{p}
  \alpha_g\bigl(S_o,\,p'\bigr)\,dp'.}
\end{equation}
При любом фиксированном $S_o$ этот интеграл определён
корректно: PVT-свойства~--- известные функции $p'$,
а ОФП при фиксированном $S_o$~--- константы.
Таким образом, \textbf{$m_g$~--- функция двух переменных
$(p, S_o)$}, определённая на всей области
$(p, S_o) \in [p_{\mathrm{ref}}, p_0] \times [0, 1]$.

\paragraph{Полный дифференциал.}
Поскольку $m_g = m_g(p, S_o)$, градиент содержит два члена:
\begin{equation}\label{eq:dmg_full}
  \boxed{\nabla m_g
  = \underbrace{\mu_{gi}B_{gi}\,\alpha_g(S_o,p)}_{
    \partial m_g/\partial p}\,\nabla p
  + \underbrace{\mu_{gi}B_{gi}\int_{p_{\mathrm{ref}}}^{p}
    \frac{\partial\alpha_g}{\partial S_o}\,dp'}_{
    \partial m_g/\partial S_o}\,\nabla S_o.}
\end{equation}
Первый член следует из теоремы Лейбница (производная
по верхнему пределу интеграла). Второй~---
производная подынтегрального выражения по параметру~$S_o$:
\begin{equation}\label{eq:dalpha_dS}
  \frac{\partial\alpha_g}{\partial S_o}
  = -\frac{k_{rg}'(S_g)}{B_g\mu_g}
  + R_s\,\frac{k_{ro}'(S_o)}{B_o\mu_o},
\end{equation}
где $k_{rg}' = dk_{rg}/dS_g$, $k_{ro}' = dk_{ro}/dS_o$.
Для степенных корреляций Кори ($k_r \sim S^n$, $n \sim 2$--$4$)
эти производные велики, и второй член
в~\eqref{eq:dmg_full} существенен.

\paragraph{Следствия для пространственного оператора.}
Подставим~\eqref{eq:dmg_full} в дивергенцию:
\begin{equation}\label{eq:div_mg_full}
  \mathrm{div}(k\nabla m_g)
  = \mathrm{div}\!\left(\frac{k\,\nabla m_g}
    {\mu_{gi}B_{gi}\alpha_g}\cdot\mu_{gi}B_{gi}\alpha_g\right)
  \neq \frac{k}{\mu_{gi}B_{gi}}\,\nabla^2 m_g.
\end{equation}
Равенство $\mathrm{div}(k\nabla m_g) \propto \nabla^2 m_g$
\textbf{не выполняется} без дополнительных условий,
поскольку $\nabla m_g$ содержит $\nabla S_o$,
а пространственный оператор в исходном
уравнении~\eqref{eq:gas_cont}~--- только $\nabla p$.
Замена $\nabla p \to \nabla m_g$ невозможна:
$\nabla m_g$ не коллинеарен $\nabla p$.

\paragraph{Замыкание $S_o(p)$: сведение к одной переменной.}
Единственный способ устранить $\nabla S_o$~--- потребовать,
чтобы $S_o$ была однозначной функцией давления:
$S_o = S_o(p)$. Тогда $\nabla S_o = (dS_o/dp)\nabla p$,
и оба слагаемых в~\eqref{eq:dmg_full} пропорциональны
$\nabla p$:
\begin{equation}\label{eq:dmg_closure}
  \nabla m_g = \frac{dm_g}{dp}\,\nabla p
  = \mu_{gi}B_{gi}\,\alpha_g\bigl(S_o(p),p\bigr)\,\nabla p.
\end{equation}
Вклад $\partial m_g/\partial S_o$ не исчез~---
он включён в полную производную $dm_g/dp$
через зависимость ОФП от $p$ посредством $S_o(p)$.
При этом пространственный оператор сводится к лапласиану:
$\mathrm{div}[k\alpha_g\nabla p]
= (1/\mu_{gi}B_{gi})\,\mathrm{div}(k\nabla m_g)
= (k/\mu_{gi}B_{gi})\nabla^2 m_g$.

\paragraph{Ограничение 1: без замыкания $m_g$ бесполезно
как замена переменной.}
Сам интеграл~\eqref{eq:mg_multi} определён корректно
при любом фиксированном~$S_o$. Проблема не в определении
$m_g$, а в его \emph{использовании} для линеаризации:
\begin{itemize}[nosep]
  \item дифференциал~\eqref{eq:dmg_full} содержит $\nabla S_o$,
        и замена $\nabla p \to \nabla m_g$ не устраняет
        нелинейность пространственного оператора;
  \item $m_g$~--- функция двух переменных:
        при одном и том же давлении $p$ разные значения
        $S_o$ дают разные~$m_g$, поэтому обращение
        $m_g \to p$ невозможно без знания~$S_o$.
\end{itemize}
Замыкание $S_o(p)$ необходимо не для определения $m_g$,
а для того, чтобы $m_g$ стало функцией одного аргумента
и замена переменной $p \to m_g$ имела смысл.

\paragraph{Замечание о Краснове.}
В документе Краснова ур.~(0.13) определяет $m_o$ через
интеграл с ОФП, а ур.~(0.14) записывает
$\nabla m_o = \alpha_o\,\mu_{oi}B_{oi}\,\nabla p$~---
без члена с $\nabla S_o$.
Замыкание $S_o(p)$ обсуждается лишь позднее
(ур.~0.33--0.36: GOR~=~const). Таким образом, ур.~(0.14)
неявно предполагает замыкание, но не оговаривает его.

\paragraph{Варианты замыканий.}
\begin{itemize}[nosep]
  \item GOR~=~const (Краснов, ур.~0.33; Boe, 1989):
        строго на квазистационарном режиме;
        приближённо при транзиенте.
  \item Формула Мартина $dS_o/dp$ при малых~$t$
        (Краснов, ур.~0.77).
  \item Численное решение уравнения Бакли--Леверетта.
\end{itemize}

\subsection{Псевдовремя $t_{ag}$: усиленная деформация}

Коэффициент при $\partial m_g/\partial t$ в квазилинейном
уравнении:
\begin{equation}\label{eq:oil_nonlin}
  \frac{A}{\alpha_g}\,\frac{\partial m_g}{\partial t}
  = \mathrm{div}(k\nabla m_g),
\end{equation}
где
\begin{align}
  \alpha_g &= \frac{k_{rg}}{B_g\mu_g}
  + R_s\frac{k_{ro}}{B_o\mu_o}, \label{eq:alpha_g}\\
  A &= \phi\left[\frac{S_g C_g}{B_g}
  + R_s'\frac{S_o}{B_o} + R_s\frac{S_o C_o}{B_o}
  + \left(\frac{R_s}{B_o} - \frac{1}{B_g}\right)
  \frac{dS_o}{dp}\right]. \label{eq:A_coeff}
\end{align}

Точное псевдовремя:
\begin{equation}\label{eq:tag_exact}
  t_{ag}(r,t) = \mu_{gi}C_{ti}\phi\int_0^t
  \frac{\alpha_g\bigl(S_o(r,\tau),\,p(r,\tau)\bigr)}
       {A\bigl(S_o(r,\tau),\,p(r,\tau)\bigr)}\,d\tau.
\end{equation}

\paragraph{Ограничение 2: двойная зависимость подинтегрального
выражения.}
В~отличие от однофазного случая, подинтегральное выражение
$\alpha_g/A$ зависит от \emph{двух} полей: $p(r,\tau)$ и
$S_o(r,\tau)$. Даже если давление осреднено
($p(r,\tau) \to \bar{p}(\tau)$), насыщенность
$S_o(r,\tau)$ остаётся пространственно неоднородной.

\paragraph{Ограничение 3: автомодельная связь $S_o(p)$.}
Осреднение $S_o(r,\tau) \to \bar{S}_o(\tau)$ предполагает
слабый пространственный контраст насыщенности.
В зоне разгазирования (ПЗП) $S_g$ может меняться
от нуля на~$r_b$ до значительных значений вблизи
скважины~--- контраст велик, и осреднение вносит
систематическую погрешность.

\paragraph{Стандартное приближение:}
\begin{equation}\label{eq:tag_std}
  t_{ag}(t) = \mu_{gi}C_{ti}\phi\int_0^t
  \frac{\alpha_g(\bar{S}_o,\bar{p})}
       {A(\bar{S}_o,\bar{p})}\,d\tau,
\end{equation}
которое корректно при:
\begin{equation}\label{eq:criterion_multi}
  \boxed{
  \left|\frac{\partial\ln(A/\alpha_g)}{\partial p}\right|
  \cdot |\delta p|
  + \left|\frac{\partial\ln(A/\alpha_g)}{\partial S_o}\right|
  \cdot |\delta S_o| \ll 1.}
\end{equation}
Здесь $\delta p = p - \bar{p}$, $\delta S_o = S_o - \bar{S}_o$.

По сравнению с однофазным критерием~\eqref{eq:criterion_1ph}
появилось \emph{второе слагаемое}: пространственная
неоднородность насыщенности. Это слагаемое может
доминировать, поскольку:
\begin{itemize}[nosep]
  \item $A/\alpha_g$ содержит ОФП $k_{ro}(S_o)$, $k_{rg}(1-S_o)$,
        которые нелинейно зависят от насыщенности
        (степенные корреляции Кори с $n \sim 2$--$4$);
  \item в ПЗП насыщенность газа меняется от $S_g = 0$
        до $S_g \sim 0{,}1$--$0{,}3$ на расстоянии
        нескольких $r_w$.
\end{itemize}

\subsection{Деформация оператора в многофазном случае}

Формула деформации~\eqref{eq:dr_ta}--\eqref{eq:laplacian_deformed}
сохраняет вид, но $\partial t_{ag}/\partial r$ определяется
зависимостью \emph{обоих} полей от~$r$:
\begin{equation}\label{eq:dtag_dr}
  \frac{\partial t_{ag}}{\partial r}
  = \mu_{gi}C_{ti}\phi\int_0^t
  \frac{\partial}{\partial r}
  \!\left[\frac{\alpha_g(S_o,p)}{A(S_o,p)}\right] d\tau.
\end{equation}
Эта производная отлична от нуля, если $p$ или $S_o$
неоднородны по~$r$. В~двухфазном случае \emph{оба} поля
неоднородны, и деформация оператора усиливается.

\subsection{Физические условия применимости}
\label{ssec:phys_conditions}

Обобщая, псевдовремя $t_{ag}$ (в стандартном приближении)
применимо при одновременном выполнении четырёх условий:

\begin{enumerate}
  \item \textbf{Слабая нелинейность $\mu_g c_g$.}
        Произведение $\mu_g c_g$ (или $A/\alpha_g$)
        слабо меняется в диапазоне $[\bar{p}-\delta p,\,
        \bar{p}+\delta p]$. Формально:
        критерий~\eqref{eq:criterion_multi}.

  \item \textbf{Слабая пространственная неоднородность
        насыщенности.}
        $|\delta S_o| \ll 1/|(\partial\ln(A/\alpha_g)/
        \partial S_o)|$. При наличии фронта разгазирования
        (скачок $S_g$ от 0 до конечных значений) это
        условие нарушается.

  \item \textbf{Корректный выбор масштаба осреднения.}
        Среднее давление $\bar{p}$ должно отражать
        состояние той области пласта, которая фактически
        <<работает>>. При BDF~--- это весь пласт.
        При транзиенте~--- область влияния
        (Anderson \& Mattar, 2005).

  \item \textbf{Существование замыкания $S_o(p)$.}
        Для вычисления $m_g(p)$ и $A(S_o,p)/\alpha_g(S_o,p)$
        нужна однозначная связь $S_o = S_o(p)$.
        Условия: GOR~=~const (квазистационар),
        или формула Мартина (малые времена),
        или автомодельное решение.
\end{enumerate}

\subsection{Сводка: однофазный vs.\ многофазный}

\begin{center}
\begin{tabular}{p{0.38\textwidth}cc}
\hline
\textbf{Характеристика} & \textbf{Однофазный газ} & \textbf{Black Oil} \\
\hline
Псевдодавление $m(p)$ & точная замена & требует $S(p)$ \\
Нелинейный коэффициент & $\mu_g c_g(p)$ & $A(S_o,p)/\alpha_g(S_o,p)$ \\
Число полей & 1 ($p$) & 2 ($p$, $S_o$) \\
Источник деформации & $\partial p/\partial r$ & $\partial p/\partial r$, $\partial S_o/\partial r$ \\
Критерий осреднения & $\delta p / p^* \ll 1$ & + $\delta S_o / S_o^* \ll 1$ \\
Замыкание & не нужно & GOR = const и~др. \\
\hline
\end{tabular}
\end{center}


%=============================================================================
\section{Неизотермический случай}\label{sec:noniso}
%=============================================================================

\subsection{Третье неизвестное поле}

При $T \neq \mathrm{const}$: PVT-свойства зависят от $(p,T)$,
в частности $R_s = R_s(p,T)$, $B_o = B_o(p,T)$,
$\mu_o = \mu_o(p,T)$. Система содержит три неизвестных поля:
$p(r,t)$, $S_o(r,t)$, $T(r,t)$.

\subsection{Дополнительные ограничения}

К ограничениям раздела~\ref{sec:multi} добавляются:

\paragraph{Ограничение 5: тройная зависимость $A/\alpha_g$.}
Коэффициент $A/\alpha_g$ зависит от трёх полей:
$p$, $S_o$, $T$. Критерий осреднения~\eqref{eq:criterion_multi}
расширяется:
\begin{equation}\label{eq:criterion_T}
  \left|\frac{\partial\ln(A/\alpha_g)}{\partial p}\right|
  |\delta p|
  + \left|\frac{\partial\ln(A/\alpha_g)}{\partial S_o}\right|
  |\delta S_o|
  + \left|\frac{\partial\ln(A/\alpha_g)}{\partial T}\right|
  |\delta T|
  \ll 1.
\end{equation}
Третье слагаемое может быть существенным из-за
чувствительности $R_s$ к температуре:
$\partial\ln R_s/\partial T \approx -4{,}5\cdot10^{-3}$~К$^{-1}$
(корреляция Стендинга).

\paragraph{Ограничение 6: смешанный градиент в $m_g$.}
Псевдодавление $m_g(p,T)$ зависит от двух переменных:
\begin{equation}
  \nabla m_g = \frac{\partial m_g}{\partial p}\nabla p
  + \frac{\partial m_g}{\partial T}\nabla T.
\end{equation}
Пространственный оператор $\mathrm{div}(k\nabla m_g)$
содержит $\nabla T$, связывая уравнение фильтрации
с уравнением энергии. Замена $p \to m_g$ \emph{не устраняет}
нелинейность пространственного оператора.

Это принципиальное отличие от однофазного случая,
где $\nabla m = m'(p)\nabla p$ и замена \emph{точна}.

\paragraph{Ограничение 7: замыкание $S_o(p,T)$.}
При $T \neq \mathrm{const}$ условие GOR~=~const
даёт $S_o = S_o(p,T)$~--- неявную зависимость
от двух переменных. Для данной точки $(r,t)$
замыкание определяет $S_o$ через \emph{локальные}
$p(r,t)$ и $T(r,t)$, но при осреднении
$T(r,\tau) \to \bar{T}(\tau)$ возникает дополнительная
погрешность: $\bar{S}_o(\tau)$ зависит от корреляции
между $p(r,\tau)$ и $T(r,\tau)$, которая теряется
при независимом осреднении.

\paragraph{Ограничение 8: уравнение энергии не линеаризуется.}
Введение $m_g$ и $t_{ag}$ не упрощает уравнение энергии.
Источники $Q_{\mathrm{JT}}$ и $Q_{\mathrm{deg}}$
содержат $\partial p/\partial t$, $\nabla p$,
$\partial S/\partial t$~--- величины, не выражающиеся
через $m_g$ без дополнительных допущений.
Система остаётся связанной.

\subsection{Полная линеаризация невозможна}

Псевдовремя $t_{ag}$ может быть формально определено,
но:
\begin{enumerate}[nosep]
  \item пространственный оператор деформирован
        (из-за $\nabla T$);
  \item уравнение энергии не поддаётся замене переменных;
  \item связь между уравнениями не разрывается.
\end{enumerate}

\textbf{Вывод:} полная линеаризация системы
фильтрация~+ энергия через псевдодавление и псевдовремя
невозможна.


%=============================================================================
\section{Стратегия выбора упрощающих допущений}\label{sec:strategy}
%=============================================================================

\subsection{Три режима по $\Delta T$}

Выбор допущения для применения псевдовремени зависит
от масштаба температурной аномалии.

\begin{center}
\renewcommand{\arraystretch}{1.3}
\begin{tabular}{p{0.15\textwidth}p{0.28\textwidth}p{0.45\textwidth}}
\hline
\textbf{Режим} & \textbf{Условие} & \textbf{Допущение} \\
\hline
Изотермический
& $\Delta T \lesssim 1$~К
& Температура не влияет на PVT.
  $m_g = m_g(p)$, $t_{ag} = t_{ag}(t)$ по Краснову.
  Стандартный подход. \\

Слабая неизотер\-мичность
& $\Delta T \sim 1$--$5$~К ($\Delta R_s/R_s \lesssim 2\%$)
& PVT-свойства линеаризуются по $\delta T$:
  $B_o(p,T) \approx B_o(p,T_0)(1+\alpha_B\,\delta T)$.
  Член $(\partial m_g/\partial T)\nabla T$~---
  поправка $O(\delta T \cdot \nabla T)$,
  пренебрежимо малая.
  Изотермические $m_g$, $t_{ag}$ применимы,
  температура~--- a~posteriori коррекция. \\

Существенная неизотер\-мичность
& $\Delta T \gtrsim 5$--$10$~К ($\Delta R_s/R_s \gtrsim 5\%$)
& Линеаризация некорректна.
  Необходима итерационная схема:
  на каждой итерации $T(r)$ фиксируется (<<замораживается>>),
  PVT вычисляются при локальном $T$, уравнение давления
  решается через $m_g(p;\,T(r))$.
  Формально: на каждой итерации~---
  изотермическая задача при неоднородной
  <<замороженной>> температуре. \\
\hline
\end{tabular}
\end{center}

\subsection{Стационар vs.\ нестационар}

\paragraph{Стационарная задача} ($\partial/\partial t = 0$).
Псевдовремя не нужно: временной оператор отсутствует.
Уравнение давления~--- уравнение Лапласа для $m_g$,
точная замена переменной. Все ограничения
разд.~\ref{ssec:conditions_1ph}--\ref{sec:noniso},
связанные с $t_{ag}$, \emph{снимаются}.

Это делает стационарную итерационную схему
(pseudopressure.tex, \S6) \textbf{свободной от ограничений
псевдовремени}. Единственные оставшиеся ограничения:
\begin{itemize}[nosep]
  \item замыкание $S_o(p)$ (GOR = const);
  \item $O(\nabla T)$-погрешность при неоднородной температуре.
\end{itemize}

\paragraph{Нестационарная задача.}
Все ограничения действуют.
Выбор допущения определяется
критерием~\eqref{eq:criterion_T}.
Дополнительная возможность: разделение масштабов
(тепловая длина $\gg$ гидродинамическая длина),
при котором температура~--- <<медленная>> переменная,
и на каждом шаге по времени можно считать $T(r)$
квази-стационарной.

\subsection{Иерархия допущений}

Упорядочим все допущения от минимального к максимальному.
Каждое следующее включает предыдущие.

\begin{enumerate}
  \item \textbf{Замыкание $S_o(p)$}~--- необходимо всегда
        (кроме однофазной зоны). Без него $m_g$
        не определено.

  \item \textbf{Осреднение $p(r,\tau) \to \bar{p}(\tau)$
        и $S_o(r,\tau) \to \bar{S}_o(\tau)$}~---
        необходимо для $t_{ag}$.
        Критерий: \eqref{eq:criterion_multi}.
        Нарушается в ПЗП при больших депрессиях.

  \item \textbf{Изотермичность / слабая неизотермичность}~---
        устраняет $\nabla T$-член в $\nabla m_g$
        и температурную зависимость $A/\alpha_g$.
        Критерий: $\Delta T \lesssim 5$~К
        (для типичных PVT-свойств).

  \item \textbf{Стационарность}~--- устраняет необходимость
        в $t_{ag}$ полностью. Наиболее надёжный подход
        для установившихся профилей.
\end{enumerate}

\paragraph{Рекомендация для ПЗП.}
В призабойной зоне пласта допущение~2
(осреднение) наиболее уязвимо: давление
падает логарифмически, и $\delta p$ в ПЗП
сравнимо с полной депрессией. Стационарная
постановка (допущение~4) позволяет обойти
эту трудность и получить установившиеся
профили $p(r)$, $S_g(r)$, $T(r)$ без
псевдовремени.

Нестационарная динамика (выход на режим,
переходные процессы) требует либо численного
решения PDE, либо метода моментов Баренблатта,
который осредняет уравнения
\emph{по пространству} (а не по времени),
избегая проблем с $t_{ag}$.

\end{document}
